\documentclass[12pt]{article}

% Packages
\usepackage[a4paper, margin=1in]{geometry} % Adjust margins
\usepackage{setspace} % For line spacing
\usepackage{tocloft} % Optional: better spacing for TOC
\usepackage[backend=biber,style=apa]{biblatex} % Load biblatex with APA style
\addbibresource{references.bib} % Link to your .bib file
\usepackage{xcolor}
\usepackage{hyperref}
\hypersetup{
	colorlinks,
	citecolor={blue!70!black},
	linkcolor={blue!70!black},
	urlcolor={blue!70!black}}
    
% Set line spacing
\onehalfspacing % or use \doublespacing for double spacing

% Title information
\title{\textbf{\large{Working Title TBD}}}
\author{\normalsize{Tim Kircali, Stefanie Stantcheva, Matthias Weber}}
\date{\normalsize{\today}}
\begin{document}
\maketitle
\pagenumbering{arabic}
\section*{Related Literature}



Our paper relates to several strands of literature. First, this study contributes to the growing body of research that measures support for economic policies using survey experiment. A part of this literature analyzes preferences over taxation and taxation-related policies. \textcite{stantcheva2021understanding} explores beliefs, reasoning and views on income taxation and estate taxation. \textcite{cruces2013biased} and \textcite{kuziemko2015elastic} analyze inequality perceptions and support of redistributive policies, \textcite{alesina2018intergenerational} and \textcite{alesina2023immigration} analyze the effects of perceptions of intergenerational mobility and immigration on preferences for redistribution. Attitudes toward tax progressivity are examined in \textcite{tarroux2019value}. \textcite{liscow2022psychology} elicit preferences over the realization rule in capital gains taxation. \textcite{fisman2020americans} explore views toward a wealth tax in the U.S. Beyond taxation, survey experiments have, for instance, been used to analyze support for legalizing payments to kidney donors \parencite{elias2019paying} or views on climate policies \parencite{dechezlepretre2025fighting}. Much of the recent research in this field builds on advances in survey methodology, which facilitate uncovering perceptions, attitudes, and reasoning underlying policy views (see \textcite{haaland2023designing} for a discussion on designing information provision experiments and \textcite{stantcheva2023run} for a guide to run survey experiments). 




Second, this project is related to the literature on tax enforcement and tax evasion, especially to the literature on the relevance of information reporting. \textcite{slemrod2007cheating} documents that misreported income in the U.S.\ varies considerably by income type. While income from wages, salaries, interest and dividends (all of which are subject to third-party reporting) is rarely misreported, income from self-employment (typically self-reported) has estimated tax gaps above $50$\%. Similarly, \textcite{kleven2011unwilling} discover near-complete compliance and no deterrence effects for third-party reported income sources in Denmark. Relatedly, \textcite{pomeranz2015no} looks at VAT paper trails in Chile and finds that the effects of randomized deterrence messages are much lower in areas covered by VAT paper trails, suggesting that evasion is much lower in these areas. Other prior work investigates effects of information mechanisms such as consumer receipt collection \parencite{naritomi2019consumers}, the role of digital payment reporting forms like 1099-K \parencite{slemrod2017does}, and discusses the theoretical foundations for why third-party reporting improves compliance \parencite{kleven2016can}. Finally, \textcite{jensen2022employment} documents how modern tax systems arise over time and finds evidence that a rise in third-party reported income drives an expansion of the income tax base. Recent research also shows that the international automatic exchange of information (AEoI) reduces offshore tax evasion \parencite{boas2024taxing,baselgia2023compliance}.\footnote{A key finding of the analysis of offshore tax evasion is that it is much more prevalent among the highest income earners. While \textcite{johns2010distribution} reveal that the net misreporting percentage generally increases with income group but peaks at the $99$th percentile, \textcite{alstadsaeter2019tax} analyze leaked data and conclude that offshore tax evasion is highly concentrated among the super-rich. They estimate that the top $0.01$\% of households evade approximately $25$\% of their taxes. By contrast, tax evasion detected in stratified random tax audits is less than $5$ percent throughout the distribution. These results are consistent with those of \textcite{boas2024taxing}, who analyze offshore tax evasion in Denmark, and \textcite{guyton2021tax}, who study tax evasion at the top of the income distribution in the U.S. They show that previous estimates had failed to capture sophisticated evasion via offshore accounts and pass-through businesses, therefore underestimating top-$1$\% tax evasion.}

Third, this project also relates to the literature on tax compliance costs and the government's role in providing services to reduce them. 
%\textcite{craig2024tax} provide a unified analysis of taxation and taxpayer education when individuals have an incomplete understanding of the tax system. 
%\textcite{feldman2016taxpayer} analyse whether tax complexity causes misperceptions about the marginal tax rate by looking at the effects of lump-sum tax liability changes. \textcite{abeler2015complex} present evidence that people underreact to changes in tax incentives when facing more complex tax systems. 
 %\textcite{blumenthal1992compliance} and \textcite{slemrod1995income} analyze the compliance costs of complex tax system for the U.S.\ individual taxpayers. 
In an early study on the tax compliance costs of individual taxpayers in the U.S., \textcite{blumenthal1992compliance} find that in 1989 taxpayers needed on average about $27$ hours and \$$66$ (1989 dollars) to file their taxes. In more recent work to understand how costly individuals perceive filing tax returns, \textcite{benzarti2020taxing} observes how taxpayers choose between itemizing deductions and claiming the standard deduction, finding that taxpayers forego large tax savings to avoid compliance costs. \textcite{hauck2024optional} shows that non-filing of tax declarations is particularly common at the lower end of the income distribution in Germany, which reduces the effectiveness of redistributive features of the tax system. \textcite{blesse2019people} present survey evidence eliciting the preferences of German taxpayers regarding the trade-off between simplifying the German tax code and reducing its progressivity. They discover that, while most taxpayers support simplifying the tax system, this support is partly driven by misconceptions about the trade-offs involved in tax complexity. \textcite{benzarti2024rising} ask U.S.\ respondents about the complexity of the tax code. They find that taxpayers perceive an increase in tax complexity and filing costs, that the vast majority of respondents consider tax filing to be tedious, and that almost all respondents would prefer a simpler tax system. Respondents would be willing to pay on average \$$130$ for a simpler tax system (despite the fact that the authors tell respondents that the simpler tax system might come with a higher tax liability due to fewer deduction possibilities). The majority of respondents is also supportive of pre-populated tax returns.

%While \textcite{zwick2021costs} explores the costs of corporate tax complexity,

%First, besides providing a more detailed measure for the support for pre-populated tax returns we also examine people’s demand for tax software. Second, by eliciting people’s views on high levels of third-party reporting and the international exchange of information, we provide novel evidence for the support of the use of tax enforcement tools involving information reporting. Third, we explore how this support interplays with support for government-provided tax filing services. For example, we investigate whether support for third-party reporting increases when taxpayers are aware that it is a prerequisite for government-provided pre-populated tax returns, or if third-party reported data were only used to reduce tax compliance costs. Finally, by eliciting important underlying considerations and beliefs, and by experimentally providing policy-relevant information, we investigate about which factual matters people might be misinformed and which considerations actually drive policy views.

Our paper contributes to the literature in several ways. We provide the first measurement of support for third-party reporting and the international exchange of tax-related information, which are of crucial importance to reducing tax evasion. Furthermore, we provide novel evidence on the support for pre-populated tax returns and government-provided free tax software. We also explore the interplay of tax evasion considerations with compliance cost considerations (such as analyzing whether the support for third-party reporting increases when taxpayers are aware that it is a prerequisite for government-provided pre-populated tax returns, or if third-party reported data were only used to reduce tax compliance costs). By eliciting underlying beliefs and by experimentally providing policy-relevant information, we shed light on which considerations drive policy views and can draw causal inferences on which information influences policy views.





\newpage
\printbibliography
\end{document}